\subsection{Phase Kickback}\label{sec:phase-kickback}

\subsubsection{Prerequisites: Eigenstates and Eigenvalues}\label{ex:prerequisites-eigenstates-and-eigenvalues}

Before introducing phase kickback, we briefly review the concepts of \textbf{eigenstates} and \textbf{eigenvalues}, which play a central role throughout quantum computing. These concepts were already introduced when explaining single-qubit gates, but they become particularly important for understanding controlled unitary operations.

\highspace
Recall that a quantum gate is represented by a \textbf{unitary operator} $U$. Most quantum states are transformed into different states when the gate is applied. However, there exists a special class of states whose direction in Hilbert space is preserved.

\highspace
\begin{flushleft}
    \textcolor{Green3}{\faIcon{bullseye} \textbf{Eigenstates}}
\end{flushleft}
A quantum state \ket{v} is called an \textbf{eigenstate} (or \textbf{eigenvector}) of a unitary operator $U$ if applying the operator changes the state only by a multiplicative constant:
\begin{equation*}
    U\ket{v} = \lambda \ket{v}
\end{equation*}
Where $\lambda$ is called the \textbf{eigenvalue} associated with the eigenstate \ket{v}.

\highspace
Unlike a generic state transformation:
\begin{equation*}
    U\ket{\psi} = \ket{\psi'}
\end{equation*}
An eigenstate is not rotated into a different direction. Instead, the operator simply multiplies the entire state by a constant factor.

\highspace
Therefore, applying $U$ to one of its eigenstates leaves the physical state unchanged except for the multiplicative coefficient.

\highspace
\begin{flushleft}
    \textcolor{Green3}{\faIcon{sync-alt} \textbf{Eigenvalues of Unitary Operators}}
\end{flushleft}
Since every quantum gate is represented by a \textbf{unitary operator}, its eigenvalues satisfy:
\begin{equation*}
    \left|\lambda\right| = 1
\end{equation*}
Consequently, every eigenvalue can always be written in the form:
\begin{equation*}
    \lambda = e^{i\phi}
\end{equation*}
Where:
\begin{itemize}
    \item $\phi$ is a real number called the \textbf{phase angle};
    \item $e^{i\phi}$ is a complex number of unit modulus.
\end{itemize}
Therefore, the eigenvalue does not change the norm of the quantum state; it only introduces a phase factor.

\highspace
The eigenvalue equation for a quantum gate is therefore usually written as:
\begin{equation*}
    U\ket{v} = e^{i\phi}\ket{v}
\end{equation*}

\highspace
\begin{flushleft}
    \textcolor{Green3}{\faIcon{balance-scale} \textbf{Special Case: Real Symmetric Unitary Operators}}
\end{flushleft}
Some quantum gates, such as the Pauli-$X$ gate, are represented by matrices that are both \textbf{real} (i.e., all entries are real numbers) and \textbf{symmetric} (i.e., the matrix equals its own transpose).

\highspace
For this particular class of operators, the eigenvalues must be real. Since unitary eigenvalues always have modulus one, the only possible values are
\begin{equation*}
    \lambda \in \left\{+1, -1\right\}
\end{equation*}
For example, the Pauli-$X$ gate has the two eigenstates:
\begin{equation*}
    X\ket{+} = +\ket{+}, \quad \text{and} \quad X\ket{-} = -\ket{-}
\end{equation*}
These two eigenstates will become extremely important when we introduce the concept of phase kickback and Deutsch's algorithm in the next sections.

\highspace
\begin{flushleft}
    \textcolor{Green3}{\faIcon{question-circle} \textbf{Why are Eigenstates important?}}
\end{flushleft}
At first sight, multiplying a state by a phase factor may seem uninteresting, since the physical state does not appear to change. However, \hl{phase kickback exploits exactly this property.}

\highspace
When a unitary operator is applied inside a \textbf{controlled operation}, the phase associated with one of its eigenstates can become a \textbf{relative phase} between components of a superposition. Unlike a global phase, a relative phase can influence interference and eventually become observable through measurement.

\highspace
For this reason, the study of phase kickback begins with the eigenvalue equation:
\begin{equation*}
    U\ket{v} = e^{i\phi}\ket{v}
\end{equation*}
Which forms the mathematical foundation of everything that follows in the next sections.